\chapter{Algebraic Bethe Ansatz and Transfer Matrices}
\label{ch:integrable-algebraic-bethe-ansatz-transfer-matrices}

\ConventionsLorentzian

Chapter~\ref{ch:integrable-yang-baxter-internal-symmetry} derived the
Yang--Baxter equation as the condition that factorized three-particle
scattering be independent of the chosen reduced decomposition.  This chapter
develops the finite-dimensional algebraic consequence needed for non-diagonal
finite-volume spectra.  The input is an \(R\)-matrix satisfying the
Yang--Baxter equation.  The output is a commuting family of transfer matrices
and, in the rational \(SU(2)\) case, explicit Bethe equations for their
simultaneous eigenvectors.

\section{Auxiliary Space, Quantum Space, and Monodromy}

Let \(V\simeq\C^d\) be a finite-dimensional complex vector space.  In spin
chains \(V\) is the one-site Hilbert space.  In finite-volume integrable QFT it
is the internal species space transported through a sequence of two-body
scatterings.  Let
\[
  R(u)\in \operatorname{End}(V\otimes V),\qquad u\in\C ,
\]
be a meromorphic endomorphism.  The parameter \(u\) is called the spectral
parameter.  For three tensor factors, \(R_{ij}(u)\) denotes \(R(u)\) acting on
the \(i\)-th and \(j\)-th factors and the identity on the remaining factor.
The Yang--Baxter equation is
\begin{equation}
  R_{12}(u-v)R_{13}(u-w)R_{23}(v-w)
  =
  R_{23}(v-w)R_{13}(u-w)R_{12}(u-v).
  \label{eq:aba-ybe-difference-form}
\end{equation}

Fix a chain length \(L\).  The quantum space is
\[
  \mathcal H_L=V_1\otimes\cdots\otimes V_L,\qquad V_j\simeq V .
\]
An additional copy \(V_0\simeq V\) is called the auxiliary space.  For complex
inhomogeneities \(\xi_1,\ldots,\xi_L\), define the monodromy matrix
\begin{equation}
  T_0(u)
  =
  R_{0L}(u-\xi_L)\cdots R_{02}(u-\xi_2)R_{01}(u-\xi_1)
  \in \operatorname{End}(V_0\otimes\mathcal H_L).
  \label{eq:aba-monodromy}
\end{equation}
The order is part of the convention.  A different order conjugates or
relabels later formulae but does not change the transfer-matrix commutativity
argument.

\begin{proposition}[RTT relation]
\label{prop:aba-rtt}
Assume \(R(u)\) satisfies \eqref{eq:aba-ybe-difference-form}.  Let \(V_0\) and
\(V_{0'}\) be two auxiliary copies of \(V\).  Then
\begin{equation}
  R_{00'}(u-v)T_0(u)T_{0'}(v)
  =
  T_{0'}(v)T_0(u)R_{00'}(u-v)
  \label{eq:aba-rtt-relation}
\end{equation}
as meromorphic endomorphisms of
\(V_0\otimes V_{0'}\otimes\mathcal H_L\).
\end{proposition}

\begin{proof}
Write
\[
  T_0(u)T_{0'}(v)
  =
  \prod_{j=L}^{1}R_{0j}(u-\xi_j)
  \prod_{j=L}^{1}R_{0'j}(v-\xi_j),
\]
where products decrease in \(j\).  The Yang--Baxter equation applied to
\(V_0\otimes V_{0'}\otimes V_j\) gives
\[
  R_{00'}(u-v)R_{0j}(u-\xi_j)R_{0'j}(v-\xi_j)
  =
  R_{0'j}(v-\xi_j)R_{0j}(u-\xi_j)R_{00'}(u-v).
\]
Operators with distinct quantum labels commute unless they share an auxiliary
factor.  Move \(R_{00'}(u-v)\) successively through the pair
\(R_{0j}(u-\xi_j)R_{0'j}(v-\xi_j)\) for \(j=L,L-1,\ldots,1\).  After the last
move all \(R_{0'j}\) factors stand to the left of all \(R_{0j}\) factors in
the order defining \(T_{0'}(v)T_0(u)\), and the auxiliary \(R_{00'}(u-v)\)
stands at the right.  This is \eqref{eq:aba-rtt-relation}.
\end{proof}

\section{Commuting Transfer Matrices}

The transfer matrix is the auxiliary trace of the monodromy:
\begin{equation}
  t(u)=\operatorname{tr}_{V_0}T_0(u)
  \in \operatorname{End}(\mathcal H_L).
  \label{eq:aba-transfer-matrix}
\end{equation}
The trace is the ordinary finite-dimensional trace on \(V_0\).  No trace is
taken over \(\mathcal H_L\).

\begin{proposition}[Transfer-matrix commutativity]
\label{prop:aba-transfer-commutativity}
Assume \(R(u)\) satisfies the Yang--Baxter equation and is invertible for
generic \(u\).  Then, for all \(u,v\) in the common domain of meromorphy,
\[
  [t(u),t(v)]=0 .
\]
\end{proposition}

\begin{proof}
Take the trace over \(V_0\otimes V_{0'}\) in the RTT relation.  On the open
set where \(R_{00'}(u-v)\) is invertible, multiply
\eqref{eq:aba-rtt-relation} on the left by \(R_{00'}(u-v)^{-1}\) and use
cyclicity of the finite-dimensional trace in the auxiliary factors:
\[
\begin{aligned}
  t(u)t(v)
  &=
  \operatorname{tr}_{0,0'}\!\left(T_0(u)T_{0'}(v)\right)  \\
  &=
  \operatorname{tr}_{0,0'}\!\left(
    R_{00'}(u-v)^{-1}T_{0'}(v)T_0(u)R_{00'}(u-v)
  \right)                                      \\
  &=
  \operatorname{tr}_{0,0'}\!\left(T_{0'}(v)T_0(u)\right)
  =
  t(v)t(u).
\end{aligned}
\]
Both sides are meromorphic in \(u,v\), hence equality on the invertibility
domain extends as a meromorphic identity.
\end{proof}

The proposition is the algebraic source of the conservation laws used in
integrable finite-volume problems.  Expanding \(t(u)\) or \(\log t(u)\) around
a regular point gives commuting operators.  Which one is the Hamiltonian is a
choice of representation and normalization, not part of the Yang--Baxter
equation alone.

\section{The Rational SU(2) R-Matrix}

Let \(V=\C^2\) with basis \(\ket{\uparrow},\ket{\downarrow}\), and let
\(P\in\operatorname{End}(V\otimes V)\) be the flip
\[
  P(x\otimes y)=y\otimes x .
\]
The rational \(SU(2)\) \(R\)-matrix in the normalization used below is
\[
  R(u)=u\,\id_{V\otimes V}+\ii P .
\]
It is \(SU(2)\)-invariant and satisfies the Yang--Baxter equation.  One direct
verification uses only
\[
  P_{ij}^2=1,\qquad
  P_{12}P_{23}P_{12}=P_{23}P_{12}P_{23},
\]
and the commutativity of disjoint flips; expanding both sides of
\eqref{eq:aba-ybe-difference-form} gives identical words in the symmetric
group algebra of \(S_3\).

For the homogeneous spin chain it is convenient to use the shifted local
operator
\[
  L_{0j}(u)=R_{0j}\!\left(u-\frac{\ii}{2}\right)
  =
  \left(u-\frac{\ii}{2}\right)\id+\ii P_{0j}.
\]
The monodromy is
\[
  T_0(u)=L_{0L}(u)\cdots L_{01}(u).
\]
Relative to the auxiliary basis, write
\[
  T_0(u)
  =
  \begin{pmatrix}
    A(u) & B(u) \\
    C(u) & D(u)
  \end{pmatrix},
\]
where \(A,B,C,D\) are operators on \(\mathcal H_L\).  The ferromagnetic
reference vector is
\[
  \Omega_L=\ket{\uparrow}^{\otimes L}.
\]
The local triangular form of \(L_{0j}(u)\) on
\(\ket{\uparrow}_j\) gives
\[
  C(u)\Omega_L=0,\qquad
  A(u)\Omega_L=\left(u+\frac{\ii}{2}\right)^L\Omega_L,\qquad
  D(u)\Omega_L=\left(u-\frac{\ii}{2}\right)^L\Omega_L .
\]

\section{Creation Operators and Bethe Equations}

The RTT relation gives commutation relations among \(A,B,C,D\).  The relations
needed to compute \(t(u)=A(u)+D(u)\) on a Bethe vector are
\begin{align}
  B(u)B(v)&=B(v)B(u),
  \label{eq:aba-bb-commutes}\\
  A(u)B(v)
  &=
  \frac{u-v+\ii}{u-v}B(v)A(u)
  -
  \frac{\ii}{u-v}B(u)A(v),
  \label{eq:aba-ab-relation}\\
  D(u)B(v)
  &=
  \frac{u-v-\ii}{u-v}B(v)D(u)
  +
  \frac{\ii}{u-v}B(u)D(v).
  \label{eq:aba-db-relation}
\end{align}
For complex numbers \(u_1,\ldots,u_M\), define
\[
  \ket{u_1,\ldots,u_M}
  =
  B(u_1)\cdots B(u_M)\Omega_L .
\]

\begin{theorem}[Algebraic Bethe ansatz for the \(XXX_{1/2}\) chain]
\label{thm:aba-xxx-bethe}
Suppose the numbers \(u_1,\ldots,u_M\) are distinct and satisfy
\begin{equation}
  \left(
    \frac{u_j+\ii/2}{u_j-\ii/2}
  \right)^L
  =
  \prod_{\substack{k=1\\ k\neq j}}^M
  \frac{u_j-u_k+\ii}{u_j-u_k-\ii},
  \qquad j=1,\ldots,M .
  \label{eq:aba-xxx-bethe-equations}
\end{equation}
Then \(\ket{u_1,\ldots,u_M}\), if nonzero, is an eigenvector of \(t(u)\) with
eigenvalue
\begin{equation}
\begin{aligned}
  \Lambda(u)
  &=
  \left(u+\frac{\ii}{2}\right)^L
  \prod_{j=1}^M\frac{u-u_j-\ii}{u-u_j}
  +
  \left(u-\frac{\ii}{2}\right)^L
  \prod_{j=1}^M\frac{u-u_j+\ii}{u-u_j}.
\end{aligned}
  \label{eq:aba-xxx-transfer-eigenvalue}
\end{equation}
\end{theorem}

\begin{proof}
Use \eqref{eq:aba-ab-relation} and \eqref{eq:aba-db-relation} to move
\(A(u)\) and \(D(u)\) through the product \(B(u_1)\cdots B(u_M)\) until they
act on \(\Omega_L\).  The terms in which \(A(u)\) or \(D(u)\) passes every
\(B(u_j)\) without changing its spectral parameter give exactly
\eqref{eq:aba-xxx-transfer-eigenvalue}.  The remaining terms have one
\(B(u_j)\) replaced by \(B(u)\).  Because the \(B\)'s commute, the coefficient
of
\[
  B(u)\prod_{\substack{k=1\\ k\neq j}}^M B(u_k)\Omega_L
\]
is proportional to
\[
  -\ii\left(u_j+\frac{\ii}{2}\right)^L
  \prod_{\substack{k=1\\ k\neq j}}^M
    \frac{u_j-u_k-\ii}{u_j-u_k}
  +
  \ii\left(u_j-\frac{\ii}{2}\right)^L
  \prod_{\substack{k=1\\ k\neq j}}^M
    \frac{u_j-u_k+\ii}{u_j-u_k}.
\]
Vanishing of this coefficient is precisely
\eqref{eq:aba-xxx-bethe-equations}.  After these unwanted terms vanish, the
displayed wanted term is the full action of \(t(u)\).
\end{proof}

The Bethe equations also have a pole-cancellation interpretation.  The
right-hand side of \eqref{eq:aba-xxx-transfer-eigenvalue} is written with
apparent simple poles at \(u=u_j\).  The residue at \(u=u_j\) vanishes exactly
when \eqref{eq:aba-xxx-bethe-equations} holds.  Thus the Bethe equations make
\(\Lambda(u)\) a polynomial of degree \(L\), as it must be because \(t(u)\) is
a polynomial operator.

\section{Hamiltonian Normalization and Momentum}

The periodic \(XXX_{1/2}\) Hamiltonian in the normalization used by the
calculation checks is
\[
  H_{XXX}=\sum_{j=1}^L(1-P_{j,j+1}),\qquad P_{L,L+1}=P_{L,1}.
\]
It is obtained from the logarithmic derivative of the transfer matrix at the
regular point, up to an additive constant fixed by
\(H_{XXX}\Omega_L=0\).  For a Bethe vector satisfying
\eqref{eq:aba-xxx-bethe-equations},
\begin{equation}
  E=\sum_{j=1}^M \frac{1}{u_j^2+1/4},
  \qquad
  \ee^{\ii P_{\rm tot}}
  =
  \prod_{j=1}^M
  \frac{u_j+\ii/2}{u_j-\ii/2}.
  \label{eq:aba-xxx-energy-momentum}
\end{equation}
For \(M=1\), set
\[
  \ee^{\ii p}=\frac{u+\ii/2}{u-\ii/2}.
\]
Then \(u=(1/2)\cot(p/2)\), the Bethe equation gives
\(\ee^{\ii pL}=1\), and
\[
  E=4\sin^2\frac p2 .
\]

\section{Use in Relativistic Integrable QFT}

In a relativistic two-dimensional integrable QFT with non-diagonal internal
scattering, the \(R\)-matrix in the monodromy is the matrix part of the
two-body \(S\)-matrix.  A particle transported once around a circle scatters
through every other particle.  The resulting internal monodromy is a transfer
matrix.  Its eigenvalues replace the scalar product of phase shifts in the
diagonal Bethe--Yang equation.  The next chapter makes this statement
explicit and then develops the nested equations that diagonalize the internal
transfer matrix for \(SU(N)\)-type scattering.
